% Part: lambda-calculus
% Chapter: introduction
% Section: fixed-point-combinator

\documentclass[../../../include/open-logic-section]{subfiles}

\begin{document}

\olfileid[id]{lam}{int}{fix}
\olsection{Kombinator Titik Tetap}

Andaikan kita mempunyai suku lambda $g$ dan menginginkan suku lain~$k$ dengan
sifat bahwa $k$ $\beta$-ekuivalen dengan $gk$. Definisikan suku-suku
\[
\fn{diag}(x) = xx
\]
dan
\[
l(x) = g(\fn{diag}(x))
\]
dengan menggunakan konvensi notasi kita; dengan kata lain, $l$ adalah suku
$\lambd[x][g(xx)]$. Misalkan $k$ adalah suku $ll$. Maka kita mempunyai
\begin{align*}
k & = (\lambd[x][g(xx)])(\lambd[x][g(xx)]) \\
& \red  g((\lambd[x][g(xx)])(\lambd[x][g(xx)])) \\
& = gk.
\end{align*}
Jika kita mengambil
\[
Y = \lambd[g][((\lambd[x][g(xx)])(\lambd[x][g(xx)]))]
\]
maka $Yg$ dan $g(Yg)$ tereduksi menjadi suatu suku yang sama; jadi
$Yg \equiv_\beta g(Yg)$. Bentuk ini dikenal sebagai ``kombinator Curry.'' Jika
sebagai gantinya kita mengambil
\[
Y = (\lambd[xg][g(xxg)])(\lambd[xg][g(xxg)])
\]
maka $Yg$ bahkan tereduksi menjadi $g(Yg)$, yang merupakan pernyataan lebih
kuat. Versi terakhir dari $Y$ ini dikenal sebagai ``kombinator Turing.''

\end{document}
